% ===== PRÉAMBULE CANONIQUE — à coller VERBATIM en tête de CHAQUE .tex =====
\documentclass[11pt,a4paper]{article}
\usepackage[utf8]{inputenc}\usepackage[T1]{fontenc}\usepackage{lmodern}
\usepackage[french]{babel}
\usepackage{amsmath,amssymb,mathtools}\usepackage{icomma}
\usepackage[version=4]{mhchem}          % \ce{...} : équations & formules
\usepackage{siunitx}\sisetup{locale=FR} % \qty{}{}, \si{}, \ang{}
\usepackage{chemfig}                    % \chemfig{...} : structures développées
\usepackage{xcolor}\usepackage{colortbl}
\usepackage{tikz}\usepackage{pgfplots}\pgfplotsset{compat=1.18}
\usepackage[most]{tcolorbox}\usepackage{enumitem}\usepackage{array,booktabs}
\usepackage[a4paper,margin=2.2cm]{geometry}
\usetikzlibrary{arrows.meta,calc,angles,quotes,positioning,patterns,backgrounds,shapes.geometric}
\definecolor{primary}{RGB}{222,49,99}\definecolor{primarydark}{RGB}{150,22,55}
\definecolor{defcol}{RGB}{0,114,178}\definecolor{methcol}{RGB}{52,92,148}
\definecolor{excol}{RGB}{0,135,90}\definecolor{attncol}{RGB}{200,40,55}
\tcbset{boxbase/.style={enhanced,breakable,boxrule=0.9pt,arc=2.5pt,
  left=3mm,right=3mm,top=2.3mm,bottom=2.3mm,fonttitle=\bfseries,coltitle=white}}
% --- boîtes TUTORIEL (noms EXACTS, ne pas renommer) ---
\newtcolorbox{cle}[1][]{boxbase,colback=defcol!6,colframe=defcol,
  title={\textsc{À retenir}\ifx\relax#1\relax\relax\else\textmd{ : \itshape #1}\fi}}
\newtcolorbox{methode}[1][]{boxbase,colback=methcol!7,colframe=methcol,sharp corners,
  title={\textsc{Méthode}\ifx\relax#1\relax\relax\else\textmd{ --- \itshape #1}\fi}}
\newtcolorbox{exemplebox}[1][]{boxbase,colback=excol!5,colframe=excol,
  title={\textsc{Exemple}\ifx\relax#1\relax\relax\else\textmd{ : \itshape #1}\fi}}
\newtcolorbox{piege}[1][]{boxbase,colback=attncol!6,colframe=attncol,
  title={\textsc{Piège}\ifx\relax#1\relax\relax\else\textmd{ : \itshape #1}\fi}}
\newtcolorbox{remarque}[1][]{boxbase,colback=black!4,colframe=black!45,
  title={\textsc{Remarque}\ifx\relax#1\relax\relax\else\textmd{ : \itshape #1}\fi}}
% --- boîtes EXERCICE / CORRIGÉ (noms EXACTS, auto-comptées) ---
\newtcolorbox[auto counter]{exo}[1][]{boxbase,colback=primary!4,colframe=primary,
  title={\textsc{Exercice}~\thetcbcounter\ifx\relax#1\relax\relax\else\textmd{ --- \itshape #1}\fi}}
\newtcolorbox[auto counter]{corrige}[1][]{boxbase,colback=excol!4,colframe=excol,
  title={\textsc{Corrigé}~\thetcbcounter\ifx\relax#1\relax\relax\else\textmd{ --- \itshape #1}\fi}}
\newcommand{\R}{\mathbb{R}}\newcommand{\N}{\mathbb{N}}\newcommand{\Z}{\mathbb{Z}}\newcommand{\ds}{\displaystyle}
% (un \usepackage{fancyhdr}+en-tête est possible ; il est ignoré côté web)
% =========================================================================
\begin{document}
\sisetup{per-mode=symbol}

\begin{exo}[qui prédomine ?]
Le couple \ce{CH3COOH/CH3COO-} a $\mathrm{p}K_a = 4{,}8$. Pour chaque pH, indique l'espèce qui prédomine.
\begin{enumerate}[label=\alph*)]
  \item $\mathrm{pH} = 2$
  \item $\mathrm{pH} = 7$
  \item $\mathrm{pH} = 4{,}8$
\end{enumerate}
\end{exo}

\begin{exo}[Henderson dans le sens direct]
Toujours pour $\mathrm{p}K_a = 4{,}8$, calcule le pH ($\mathrm{pH} = \mathrm{p}K_a + \log\frac{[\ce{A-}]}{[\ce{HA}]}$).
\begin{enumerate}[label=\alph*)]
  \item $\dfrac{[\ce{A-}]}{[\ce{HA}]} = 1$
  \item $\dfrac{[\ce{A-}]}{[\ce{HA}]} = 10$
  \item $\dfrac{[\ce{A-}]}{[\ce{HA}]} = 0{,}1$
\end{enumerate}
\end{exo}

\begin{exo}[$\Delta\mathrm{p}K_a$ et constante d'équilibre]
Calcule $\Delta\mathrm{p}K_a = \mathrm{p}K_a(\text{formé}) - \mathrm{p}K_a(\text{réactif})$ puis
$K_{\text{équi}} = 10^{\Delta\mathrm{p}K_a}$.
\begin{enumerate}[label=\alph*)]
  \item acide réactif $\mathrm{p}K_a = 2{,}1$ ; acide formé $\mathrm{p}K_a = 10{,}4$
  \item acide réactif $\mathrm{p}K_a = 6{,}4$ ; acide formé $\mathrm{p}K_a = 4{,}8$
\end{enumerate}
\end{exo}

\begin{exo}[complète, équilibrée ou impossible ?]
À partir de $\Delta\mathrm{p}K_a$, classe chaque réaction.
\begin{enumerate}[label=\alph*)]
  \item $\Delta\mathrm{p}K_a = 5$
  \item $\Delta\mathrm{p}K_a = 1$
  \item $\Delta\mathrm{p}K_a = -4$
\end{enumerate}
\end{exo}

\begin{exo}[réagir avec sa propre base conjuguée ?]
On envisage de faire réagir l'acide acétique \ce{CH3COOH} ($\mathrm{p}K_a = 4{,}8$) avec l'ion acétate
\ce{CH3COO-}.
\begin{enumerate}[label=\alph*)]
  \item Quel serait l'« acide formé » ? En déduis $\Delta\mathrm{p}K_a$ et $K_{\text{équi}}$.
  \item La réaction a-t-elle lieu dans un sens privilégié ?
\end{enumerate}
\end{exo}

\end{document}
